%% Question bank chapter file - generated by bank_tools.py build (do not edit while the bank is still being built:
%% it is regenerated from the per-test data in Question Bank\data). Every question is one bankq block:
%%   \begin{bankq}[id=..., source={...}, paper=.., qno=.., type=..., subtopic={...}, answer={...}, answer_src=key|solved,
%%                 working={...}, difficulty=1-10, diff_parts={K S I C F}, diff_note={...},
%%                 group=<passage id or empty>, also_in={duplicate copies in other tests}]
%% difficulty = min(10, K+S+I+C+F-1): K concepts 1-3, S steps 1-4, I insight 0-2, C computation 0-2, F format 0-1
%%   ...question LaTeX, options, figures (\bfig / \bfigi -> img/)...
%%   \end{bankq}
%% Passages shared by several questions are bankpassage blocks; their questions carry group=<passage id>.
\documentclass[10pt]{article}
\usepackage[a4paper,margin=1.8cm]{geometry}
\usepackage[T1]{fontenc}
\usepackage{amsmath,amssymb,mathtools}
\usepackage{graphicx}
\usepackage[export]{adjustbox}
\usepackage{xcolor,needspace,keyval}
\usepackage{enumitem}
\usepackage{booktabs,array,tabularx}
\usepackage[hidelinks]{hyperref}
\setlength{\parindent}{0pt}
\setlength{\parskip}{3pt}
\makeatletter
\define@key{bq}{id}{\def\bq@id{#1}}\define@key{bq}{source}{\def\bq@src{#1}}\define@key{bq}{paper}{\def\bq@paper{#1}}
\define@key{bq}{qno}{\def\bq@qno{#1}}\define@key{bq}{type}{\def\bq@type{#1}}\define@key{bq}{subtopic}{\def\bq@sub{#1}}
\define@key{bq}{answer}{\def\bq@ans{#1}}\define@key{bq}{answer_src}{\def\bq@asrc{#1}}\define@key{bq}{working}{\def\bq@work{#1}}
\define@key{bq}{group}{\def\bq@group{#1}}\define@key{bq}{also_in}{\def\bq@also{#1}}
\define@key{bq}{difficulty}{\def\bq@diff{#1}}\define@key{bq}{diff_parts}{\def\bq@dparts{#1}}\define@key{bq}{diff_note}{\def\bq@dnote{#1}}
\newenvironment{bankq}[1][]{\def\bq@work{}\def\bq@also{}\def\bq@group{}\def\bq@diff{?}\def\bq@dparts{}\def\bq@dnote{}\setkeys{bq}{#1}%
  \par\needspace{4\baselineskip}\medskip\noindent{\small\textbf{\bq@id}\hfill\color{gray}[\bq@src, Paper \bq@paper, Q\bq@qno]
  \bq@type\ $\cdot$ \bq@sub\ $\cdot$ \textbf{D\bq@diff}}\par\nopagebreak\relax}
  {\par{\small\color{teal}Answer: \textbf{\bq@ans}\ifx\bq@work\@empty\else\ (solved: \bq@work)\fi
  \ifx\bq@also\@empty\else\par\color{gray}Also in: \bq@also\fi}\par}
\define@key{bp}{id}{\def\bp@id{#1}}\define@key{bp}{source}{\def\bp@src{#1}}\define@key{bp}{questions}{\def\bp@qs{#1}}
\newenvironment{bankpassage}[1][]{\setkeys{bp}{#1}\par\needspace{4\baselineskip}\bigskip\noindent
  {\small\textbf{Passage \bp@id}\hfill\color{gray}[\bp@src; questions \bp@qs]}\par\nopagebreak\relax\itshape}{\par}
\makeatother
\newcommand{\bfig}[2]{\par\smallskip{\centering\includegraphics[width=#1pt,max width=0.62\linewidth,max height=0.32\textheight]{img/#2}\par}\smallskip}
\newcommand{\bfigi}[2]{\includegraphics[width=#1pt,max width=\linewidth,max height=0.32\textheight,valign=c]{img/#2}}
\newcommand{\bmissing}[1]{\textit{[figure missing in the source conversion: \detokenize{#1}]}}
\newenvironment{qoptions}{\begin{itemize}[leftmargin=1.8em,itemsep=1pt,topsep=2pt,label={}]}{\end{itemize}}
\newcommand{\opt}[2]{\item[(#1)] #2}
\providecommand{\fwstart}{}\providecommand{\fwend}{}\providecommand{\qfullsplit}{}
\newcolumntype{L}[1]{>{\raggedright\arraybackslash}p{#1}}
\newcolumntype{C}[1]{>{\centering\arraybackslash}p{#1}}
\newcolumntype{R}[1]{>{\raggedleft\arraybackslash}p{#1}}
\providecommand{\sectionhdr}[1]{}\providecommand{\subsectionhdr}[1]{}


\begin{document}
{\Large\bfseries Waves --- Question Bank}\par\smallskip
{\small Physics chapter P18. Questions from all converted test papers, grouped by source series. Answers: printed key where it matches the paper, otherwise solved (answer\_src=solved, with working). Exact repeats across papers appear once, with also\_in listing the other copies.}\par\medskip\hrule

\section*{Score-I papers}
\begin{bankq}[id=SI12-P1-PHY-S1-Q09, source={Score-I - Test 12 (09 Mar 2022)}, paper=1, qno=9, type=multi, subtopic={Doppler effect}, answer={A,B}, answer_src=key, difficulty=2, diff_parts={K1 S1 I0 C0 F1}, diff_note={concepts: Doppler effect with wind | steps: effective sound speed relative to the medium for each wind direction; frequency formula and comparison with f1}]
Two vehicles, each moving with speed $u$ on the same horizontal straight road, are approaching each other. Wind blows along the road with velocity $w$. One of these vehicles blows a whistle of frequency $f_1$. An observer in the other vehicle hears the frequency of the whistle to be $f_2$. The speed of sound in still air is $V$. The correct statement(s) is (are):
\begin{qoptions}
\opt{A}{If the wind blows from the observer to the source, $f_2>f_1$}
\opt{B}{If the wind blows from the source to the observer, $f_2>f_1$}
\opt{C}{If the wind blows from observer to the source, $f_2<f_1$}
\opt{D}{If the wind blows from the source to the observer, $f_2<f_1$}
\end{qoptions}
\end{bankq}
\begin{bankq}[id=SI12-P1-PHY-S2-Q04, source={Score-I - Test 12 (09 Mar 2022)}, paper=1, qno=4, type=numerical, subtopic={Standing waves}, answer={2.00}, answer_src=key, difficulty=4, diff_parts={K1 S2 I1 C0 F1}, diff_note={concepts: energy in a standing wave | steps: mode number and time dependence of the phase; energy fully potential when cos(wt + pi/3) = +-1; second such instant gives 120 t0 | I: energy is purely potential at maximum displacement}]
A string fixed at both the ends of length 2 m, vibrating in its 7th overtone. Equation of the standing wave is given by $y=A\sin kx\cos(\omega t+\pi/3)$. All the symbols have their usual meaning. Mass per unit length of the string is 0.5 gm/cm. Given that $A=1$ cm and $\omega=100\pi$ rad/sec. If after time $t_0$ starting from $t=0$ sec, energy of vibration is completely potential for the second time. Find $120\,t_0$. (Use $\pi^2=10$)
\end{bankq}

\section*{TOAS papers}
\begin{bankq}[id=TOAS15-P1-PHY-S2-Q08, source={TOAS - Test 15 (22 Aug 2022)}, paper=1, qno=8, type=multi, subtopic={Waves on a string}, answer={A,C,D}, answer_src=key, difficulty=7, diff_parts={K2 S3 I1 C1 F1}, diff_note={concepts: travelling wave equation, energy flow on a string | steps: amplitude and angular frequency from the figure; wave number and direction from the slope and velocity of P; phase; tension from v\^{}2 mu; energy flow rate}]
The figure shows a snap shot of a vibrating string of linear mass density 0.1 kg/m at $t=0$. The particle P is observed moving up with velocity $20\sqrt3$cm/s. The tangent at P makes an angle $60^\circ$ with x-axis then
\bfig{166.1}{TOAS15_p1_phy_s2_q08_1.png}
\begin{qoptions}
\opt{A}{The wave equation is given by $y=4\text{mm}\sin\left[10\pi t+50\pi x+\dfrac\pi4\right]$.}
\opt{B}{The wave equation is given by $y=4\text{mm}\cos\left[10\pi t-50\pi x+\dfrac\pi4\right]$.}
\opt{C}{Rate of flow of energy is towards left and it's value at point P is $12\times10^{-5}$J}
\opt{D}{Tension in string is $4\times10^{-3}$ N.}
\end{qoptions}
\end{bankq}
\begin{bankq}[id=TOAS15-P1-PHY-S2-Q09, source={TOAS - Test 15 (22 Aug 2022)}, paper=1, qno=9, type=multi, subtopic={Standing waves}, answer={A,B,C,D}, answer_src=key, difficulty=4, diff_parts={K2 S2 I0 C0 F1}, diff_note={concepts: waves on strings of different density, standing waves | steps: speeds from tension and density; condition for node or antinode at P; frequencies of the first modes}]
Two string having linear mass density 0.1 kg/m and 0.4kg/m are connected as shown and pulled by force 10 N. Select correct alternative :-
\bfig{132.2}{TOAS15_p1_phy_s2_q09_1.png}
\begin{qoptions}
\opt{A}{Fundamental frequency of system is 10Hz, if antinote is formed at point P.}
\opt{B}{Fundamental frequency of system is 5Hz, if node is formed at point P.}
\opt{C}{$1^{\text{st}}$ overtone has frequency 30Hz, if antinote is formed at point P.}
\opt{D}{$1^{\text{st}}$ overtone has frequency 10Hz, if node is formed at point P.}
\end{qoptions}
\end{bankq}
\begin{bankq}[id=TOAS15-P2-PHY-S1-Q02, source={TOAS - Test 15 (22 Aug 2022)}, paper=2, qno=2, type=multi, subtopic={Waves on a string}, answer={B,C,D}, answer_src=key, difficulty=3, diff_parts={K1 S2 I0 C0 F1}, diff_note={concepts: travelling pulse | steps: direction from the sign of 4x + 5t; speed 5/4; maximum displacement; symmetry}]
$Y(x,t)=\dfrac{0.8}{\left[(4x+5t)^2+5\right]}$ represents a moving pulse where $x$ and $y$ are in metres and $t$ in second. Then :
\begin{qoptions}
\opt{A}{pulse is moving in positive x-direction.}
\opt{B}{in 2s it will travel a distance of 2.5 m.}
\opt{C}{maximum displacement of particles is 0.16 m.}
\opt{D}{it is a symmetric pulse at $t=0$.}
\end{qoptions}
\end{bankq}
\begin{bankq}[id=TOAS16-P1-PHY-S1-Q03, source={TOAS - Test 16 (04 Sep 2022)}, paper=1, qno=3, type=single, subtopic={Waves on a string}, answer={D}, answer_src=key, difficulty=1, diff_parts={K1 S1 I0 C0 F0}, diff_note={concepts: transverse velocity in a travelling wave | steps: v = omega sqrt(A\^{}2 - y\^{}2) with the given displacement}]
The function $y(x,t)=(15.0\text{ cm})\cos(\pi x-15\pi t)$, with $x$ in meters and $t$ in seconds, describes a wave on a taut string. What is the transverse speed for a point on the string at an instant when that point has the displacement $y=+12.0$ cm ?
\begin{qoptions}
\opt{A}{$180\pi$ cm/s}
\opt{B}{$225\pi$ cm/s}
\opt{C}{$45\pi$ cm/s}
\opt{D}{$135\pi$ cm/s}
\end{qoptions}
\end{bankq}
\begin{bankq}[id=TOAS16-P1-PHY-S2-Q08, source={TOAS - Test 16 (04 Sep 2022)}, paper=1, qno=8, type=multi, subtopic={Standing waves}, answer={B,C,D}, answer_src=key, difficulty=4, diff_parts={K2 S2 I0 C0 F1}, diff_note={concepts: standing waves on a rope, wave speed | steps: k and omega from the equation; length from the second harmonic; mass per length from v\^{}2 = T/mu; phases of the loops}]
A rope, under a tension of 200 N and fixed at both ends, oscillates in a second-harmonic standing wave pattern. The displacement of the rope is given by $y=(0.10\text{m})(\sin\pi x/2)\sin12\pi t$, where $x=0$ at one end of the rope, $x$ is in meters, and $t$ is in seconds.
\begin{qoptions}
\opt{A}{The length of the rope is 6m.}
\opt{B}{Mass/length of the rope is $\dfrac{25}{72}$ kg/m .}
\opt{C}{There is an antinode at $x=3$m.}
\opt{D}{The particle at $x=1$m \& $x=1.5$ m always oscillate in same phase.}
\end{qoptions}
\end{bankq}
\begin{bankq}[id=TOAS16-P1-PHY-S2-Q09, source={TOAS - Test 16 (04 Sep 2022)}, paper=1, qno=9, type=multi, subtopic={Waves on a string}, answer={B,D}, answer_src=key, difficulty=6, diff_parts={K2 S2 I1 C1 F1}, diff_note={concepts: wave speed on a hanging rope, travel time | steps: tension at A and B; wavelength proportional to speed; integrate dx/v with tension varying linearly}]
A block of mass $3m$ is suspended vertically by a rope of mass $m$ and length $\ell$ as shown in figure. A wave pulse of wavelength $\lambda_0$ is generated at bottom most point A of the string.
\bfig{108.3}{TOAS16_p1_phy_s2_q09_1.png}
\begin{qoptions}
\opt{A}{Wavelength of wave pulse at point B is $\sqrt{\dfrac34}\,\lambda_0$.}
\opt{B}{Wavelength of wave pulse at point B is $\sqrt{\dfrac43}\,\lambda_0$.}
\opt{C}{Time taken by wave pulse to reach point B is $\left(2-\sqrt3\right)\sqrt{\dfrac\ell g}$}
\opt{D}{Time taken by wave pulse to reach point B is $2\left(2-\sqrt3\right)\sqrt{\dfrac\ell g}$.}
\end{qoptions}
\end{bankq}
\begin{bankq}[id=TOAS16-P1-PHY-S4-Q07, source={TOAS - Test 16 (04 Sep 2022)}, paper=1, qno=7, type=numerical, subtopic={Sound waves}, answer={40}, answer_src=key, difficulty=2, diff_parts={K1 S1 I0 C0 F1}, diff_note={concepts: intensity and loudness level | steps: I = P/4 pi r\^{}2; beta = 10 log(I/I0)}]
A point sound source has an output of $4\pi\mu$W. What is loudness level in (dB) at a distance of 10 m from the source ?
\end{bankq}
\begin{bankq}[id=TOAS16-P2-PHY-S1-Q01, source={TOAS - Test 16 (04 Sep 2022)}, paper=2, qno=1, type=multi, subtopic={Waves on a string}, answer={B,C}, answer_src=key, difficulty=3, diff_parts={K1 S1 I1 C0 F1}, diff_note={concepts: wave train from the waveform | steps: spacing of crests gives frequency change; speed of a particle from the slope}]
A boy is jerking one end of a uniform taut string. The wave train propagating to the right has been shown in the figure.
\bfig{160.9}{TOAS16_p2_phy_s1_q01_1.png}
\begin{qoptions}
\opt{A}{The boy is jerking at gradually decreasing frequency.}
\opt{B}{The boy is jerking at gradually increasing frequency.}
\opt{C}{Speed of particle at a is more than speed of particle at b.}
\opt{D}{Speed of particle at b is more than speed of particle at a.}
\end{qoptions}
\end{bankq}
\begin{bankpassage}[id=TOAS16-P2-PHY-S2-G9-10, source={TOAS - Test 16 (04 Sep 2022)}, questions={9-10}]
In resonance column experiment a tuning fork of frequency $f=400$ Hz is held above the pipe as shown in figure. The reservoir is raised and lowered to change the level of water and thus the length of the column of air in the tube. The area of cross section of the reservoir is 6 times that of the pipe. Initially, the reservoir is kept so that the pipe is full up to the brim. Tuning fork is sounded and the reservoir is lowred. When the reservoir is lowred by 21 cm, first resonance is recorded. When the reservoir is lowered further by 49 cm the second resonance is heard.
\bfig{82.1}{TOAS16_p2_phy_s2_para9_10_1.png}
\end{bankpassage}
\begin{bankq}[id=TOAS16-P2-PHY-S2-Q09, source={TOAS - Test 16 (04 Sep 2022)}, paper=2, qno=9, type=single, subtopic={Standing waves}, answer={D}, answer_src=key, difficulty=2, diff_parts={K1 S1 I1 C0 F0}, diff_note={concepts: resonance column with a reservoir | steps: pipe level moves 6/7 of the reservoir displacement; area from the geometry}, group=TOAS16-P2-PHY-S2-G9-10]
The area of reservoir is :-
\begin{qoptions}
\opt{A}{$54\pi$ cm$^2$}
\opt{B}{$112\pi$ cm$^2$}
\opt{C}{$27\pi$ cm$^2$}
\opt{D}{$150\pi$ cm$^2$}
\end{qoptions}
\end{bankq}
\begin{bankq}[id=TOAS16-P2-PHY-S2-Q10, source={TOAS - Test 16 (04 Sep 2022)}, paper=2, qno=10, type=single, subtopic={Standing waves}, answer={B}, answer_src=key, difficulty=3, diff_parts={K1 S2 I1 C0 F0}, diff_note={concepts: resonance column | steps: pipe level change between resonances; half wavelength; v = f lambda}, group=TOAS16-P2-PHY-S2-G9-10]
The speed of sound is :-
\begin{qoptions}
\opt{A}{228 m/s}
\opt{B}{336 m/s}
\opt{C}{398 m/s}
\opt{D}{328 m/s}
\end{qoptions}
\end{bankq}
\begin{bankq}[id=TOAS16-P2-PHY-S4-Q08, source={TOAS - Test 16 (04 Sep 2022)}, paper=2, qno=8, type=numerical, subtopic={Waves on a string}, answer={6}, answer_src=key, difficulty=4, diff_parts={K1 S2 I1 C0 F1}, diff_note={concepts: velocity graph of a travelling wave | steps: amplitude of velocity from the graph; omega; acceleration at the point from the phase}]
A travelling wave is propagating along a tight rope. Velocity of any point of rope is shown in the graph. There is a point $P$ on rope. If frequency of wave is $\dfrac{\sqrt3}4$Hz then magnitude of acceleration of $P$ is ?
\bfig{119.0}{TOAS16_p2_phy_s4_q08_1.png}
\end{bankq}
\begin{bankq}[id=TOAS16-P2-PHY-S4-Q09, source={TOAS - Test 16 (04 Sep 2022)}, paper=2, qno=9, type=numerical, subtopic={Standing waves}, answer={4}, answer_src=key, difficulty=3, diff_parts={K1 S2 I0 C0 F1}, diff_note={concepts: energy of a standing wave | steps: speed on the string; frequency of the third overtone; energy of the oscillation}]
A string of mass 1gm, length 3m, tension 30 N is fixed at both the ends and is oscillating in 3rd overtone. The amplitude at the antinode is 0.1 mm. What is the energy of oscillations (in $\mu$J) ?
\end{bankq}
\begin{bankq}[id=TOAS17-P1-PHY-S1-Q01, source={TOAS - Test 17 (18 Sep 2022)}, paper=1, qno=1, type=multi, subtopic={Waves on a string}, answer={A,B,C,D}, answer_src=key, difficulty=7, diff_parts={K2 S3 I1 C1 F1}, diff_note={concepts: waves on a rotating rope, centripetal tension | steps: tension at distance r from the centripetal condition; wave speed and wavelength along the rope; travel time by integration; speed in the ground frame; amplitude change}]
A block of mass $M$ is connected to a uniform rope of length $L$ having uniform mass per unit length $\mu$. The other end of rope is hinged at $O$ and system is rotating with constant angular velocity $\omega$ in gravity free space. A transverse pulse of wavelength $\lambda$ is produced at $O$. Assume no loss in energy. Then (Given : $\mu=\dfrac{2M}{3L}$)
\bfig{161.1}{TOAS17_p1_phy_s1_q01_1.png}
\begin{qoptions}
\opt{A}{Wavelength of pulse decrease as it reaches point $A$ and becomes $\dfrac{\sqrt3}{2}\lambda$.}
\opt{B}{Time taken by pulse to travel from $O$ to $A$ is $\dfrac{\pi}{3\sqrt2\,\omega}$.}
\opt{C}{As pulse propagate from $O$ to $A$, it's Amplitude continuously increases.}
\opt{D}{Pulse speed at point $A$ in frame of ground is $\sqrt{\dfrac52}\,\omega L$.}
\end{qoptions}
\end{bankq}
\begin{bankq}[id=TOAS17-P1-PHY-S1-Q05, source={TOAS - Test 17 (18 Sep 2022)}, paper=1, qno=5, type=multi, subtopic={Standing waves}, answer={A,B,C,D}, answer_src=key, difficulty=5, diff_parts={K2 S2 I0 C1 F1}, diff_note={concepts: wave speed in a rod and a gas, standing waves | steps: speed in the rod from Y/rho; speed in the gas from B/rho; frequency from node spacing; mode of the rod from its clamping}]
A rod of length $L=1$m is made of material having density $\rho=\dfrac{3g}{cc}$ and youngs modulus $Y=2.43\times10^9$ N/m$^2$, is clamped at distance $\dfrac L4$ from one end and other end has as metal disk attached. Glass tubs contain gas of density $\rho_0=0.01$ g/cc and bulk modulus $B_0=9\times10^5$ N/m$^2$. When longitudanal oscillation are setup in rod then fine powder present in tube pill up at regular interval of $\ell=3.33$ cm, then select correct alternatives.
\bfig{266.7}{TOAS17_p1_phy_s1_q05_1.png}

\begin{qoptions}
\opt{A}{Speed of wave in rod is 900m/s.}
\opt{B}{Speed of sound in gas 300 m/s}
\opt{C}{frequency of wave is nearly 4500 Hz.}
\opt{D}{Rod is vibrating in 2$^{\text{nd}}$ overtone.}
\end{qoptions}
\end{bankq}
\begin{bankq}[id=TOAS17-P1-PHY-S3-Q02, source={TOAS - Test 17 (18 Sep 2022)}, paper=1, qno=2, type=numerical, subtopic={Doppler effect}, answer={4.12 to 4.15}, answer_src=key, difficulty=6, diff_parts={K1 S2 I2 C1 F1}, diff_note={concepts: Doppler effect with wind and a moving reflector | steps: velocities relative to the air; frequency at the truck; reflected frequency at the car; beats with the direct sound}]
Wind is blowing due east with speed 10 m/s. A car is moving west with speed 20 m/s towards a truck moving towards car with speed 10 m/s. If car blow horn of frequency 20 Hz, then find number of beats heard per second by a person in car. Assume waves reflected from truck reaches car. (speed of sound in air = 320 m/s)
\end{bankq}
\begin{bankq}[id=TOAS17-P2-PHY-S1-Q03, source={TOAS - Test 17 (18 Sep 2022)}, paper=2, qno=3, type=single, subtopic={Standing waves}, answer={A}, answer_src=key, difficulty=1, diff_parts={K1 S1 I0 C0 F0}, diff_note={concepts: vibrations of an open pipe | steps: harmonics of an open pipe; pressure node and displacement antinode at the open end}]
If we study the vibration of a pipe open at both ends, then the following statements is not true :
\begin{qoptions}
\opt{A}{Only odd harmonics of the fundamental frequency will be generated.}
\opt{B}{All harmonics of the fundamental frequency will be generated.}
\opt{C}{Open end will be pressure node}
\opt{D}{Open end will be displacement antinode}
\end{qoptions}
\end{bankq}
\begin{bankq}[id=TOAS17-P2-PHY-S1-Q04, source={TOAS - Test 17 (18 Sep 2022)}, paper=2, qno=4, type=single, subtopic={Standing waves}, answer={A}, answer_src=key, difficulty=1, diff_parts={K1 S1 I0 C0 F0}, diff_note={concepts: standing wave between nodes | steps: distance between end nodes and wavelength}]
A standing wave having 3 nodes and 2 antinodes is formed between two atoms having a distance of 1.21 \AA\ between them. The wavelength of the standing wave is -
\begin{qoptions}
\opt{A}{1.21 \AA}
\opt{B}{2.42 \AA}
\opt{C}{3.63 \AA}
\opt{D}{6.05 \AA}
\end{qoptions}
\end{bankq}
\begin{bankq}[id=TOAS17-P2-PHY-S1-Q05, source={TOAS - Test 17 (18 Sep 2022)}, paper=2, qno=5, type=single, subtopic={Standing waves}, answer={A}, answer_src=key, difficulty=1, diff_parts={K1 S1 I0 C0 F0}, diff_note={concepts: closed and open pipes | steps: quarter and half wavelength for the same fundamental}]
A closed organ pipe and an open organ pipe are tuned to the same fundamental frequency. What is the ratio of their lengths ?
\begin{qoptions}
\opt{A}{1 : 2}
\opt{B}{2 : 1}
\opt{C}{2 : 3}
\opt{D}{4 : 3}
\end{qoptions}
\end{bankq}
\begin{bankq}[id=TOAS17-P2-PHY-S1-Q06, source={TOAS - Test 17 (18 Sep 2022)}, paper=2, qno=6, type=single, subtopic={Doppler effect}, answer={A}, answer_src=key, difficulty=1, diff_parts={K1 S1 I0 C0 F0}, diff_note={concepts: Doppler effect with a moving observer | steps: f = f0 (v + v0)/v}]
A traveller is running with velocity $v_0$ towards a stationary train. If the train sounds a horn of frequency $n$ then the apparent frequency heard by the traveller will be ($v$ is speed of sound) :
\begin{qoptions}
\opt{A}{$n'=n\left(\dfrac{v+v_0}{v}\right)$}
\opt{B}{$n'=n\left(\dfrac{v-v_0}{v+v_s}\right)$}
\opt{C}{$n'=n\left(\dfrac{v-v_0}{v}\right)$}
\opt{D}{None of these}
\end{qoptions}
\end{bankq}
\begin{bankq}[id=TOAS17-P2-PHY-S1-Q07, source={TOAS - Test 17 (18 Sep 2022)}, paper=2, qno=7, type=single, subtopic={Waves on a string}, answer={D}, answer_src=key, difficulty=1, diff_parts={K1 S1 I0 C0 F0}, diff_note={concepts: travelling wave equation | steps: speed = frequency times wavelength from the equation}]
The equation of a progressive wave is $y=0.02\sin2\pi\left[\dfrac t{0.01}-\dfrac x{0.3}\right]$, where $x$ and $y$ are in meters and $t$ is in second. The velocity (in m/s) of propagation of the wave is :
\begin{qoptions}
\opt{A}{400}
\opt{B}{40}
\opt{C}{300}
\opt{D}{30}
\end{qoptions}
\end{bankq}
\begin{bankq}[id=TOAS17-P2-PHY-S1-Q08, source={TOAS - Test 17 (18 Sep 2022)}, paper=2, qno=8, type=single, subtopic={Waves on a string}, answer={A}, answer_src=key, difficulty=1, diff_parts={K1 S1 I0 C0 F0}, diff_note={concepts: period of a wave | steps: quarter period is 0.17 s; frequency}]
In a sinusoidal wave, the time required for a particular point to move from maximum displacement to zero displacement is 0.17 sec. The frequency of the wave is :-
\begin{qoptions}
\opt{A}{1.47 Hz}
\opt{B}{0.36 Hz}
\opt{C}{0.73 Hz}
\opt{D}{2.94 Hz}
\end{qoptions}
\end{bankq}
\begin{bankq}[id=TOAS17-P2-PHY-S1-Q13, source={TOAS - Test 17 (18 Sep 2022)}, paper=2, qno=13, type=single, subtopic={Waves on a string}, answer={A}, answer_src=key, difficulty=3, diff_parts={K1 S2 I1 C0 F0}, diff_note={concepts: wave speed on a non-uniform rope | steps: tension from the weight below; mass per length; speed squared against position}]
A non-uniform rope of length $\ell$ hangs from a ceilling. Mass per unit length of rope ($\mu$) changes as $\mu=\mu_0e^y$, where $y$ is the distance along the string from its lowest point. Then graph between square of velocity of wave $\left(v_y^2\right)$ and $y$ will be best represented as :
\bfig{42.4}{TOAS17_p2_phy_s1_q13_1.png}
\begin{qoptions}
\opt{A}{\bfigi{81.1}{TOAS17_p2_phy_s1_q13_A.png}}
\opt{B}{\bfigi{81.1}{TOAS17_p2_phy_s1_q13_B.png}}
\opt{C}{\bfigi{89.1}{TOAS17_p2_phy_s1_q13_C.png}}
\opt{D}{None of these}
\end{qoptions}
\end{bankq}
\begin{bankq}[id=TOAS17-P2-PHY-S2-Q05, source={TOAS - Test 17 (18 Sep 2022)}, paper=2, qno=5, type=numerical, subtopic={Beats}, answer={100}, answer_src=key, difficulty=2, diff_parts={K1 S1 I0 C0 F1}, diff_note={concepts: beat frequency | steps: frequencies of the two forks; difference}]
The frequency of a fork A is 3\% more than the frequency of a standard fork whereas the frequency of fork B is 3\% less. The forks A and B produce 6 beats per second. The frequency of standard fork (in Hz) will be :-
\end{bankq}
\begin{bankq}[id=TOAS17-P2-PHY-S2-Q10, source={TOAS - Test 17 (18 Sep 2022)}, paper=2, qno=10, type=numerical, subtopic={Sound waves}, answer={100}, answer_src=key, difficulty=2, diff_parts={K1 S1 I0 C0 F1}, diff_note={concepts: decibel scale | steps: 20 dB means a factor of 100}]
A sound absorber attenuates the sound level by 20 dB. The intensity decreases by a factor of :
\end{bankq}
\begin{bankq}[id=TOAS18-P1-PHY-S2-Q10, source={TOAS - Test 18 (30 Oct 2022)}, paper=1, qno=10, type=single, subtopic={Waves on a string}, answer={A}, answer_src=solved, working={I: $\lambda=40$ cm, $f=50$ Hz, $v_p=\omega A|\cos|=\pi/10$ with the particle moving up; II and IV standing waves have unknown phase of time, so the direction is undetermined (T); III longitudinal wave moving to $-x$ with $v_p=\pi/10$}, difficulty=5, diff_parts={K2 S2 I1 C1 F0}, diff_note={concepts: particle velocity in travelling, standing waves | steps: wavelength and frequency from the graph; velocity amplitude at the point; direction from the slope; the undetermined standing wave cases}]
List-I contain information about mechanical waves and list-II contain information about velocity of particle of medium corresponding to point $P$ shown. Match list-I with list-II using given codes.
\begin{center}\footnotesize\setlength{\tabcolsep}{2pt}
\begin{tabular}{|L{0.7cm}L{3.5cm}|L{0.6cm}L{2.2cm}|}\hline
\multicolumn{2}{|c|}{List-I}&\multicolumn{2}{c|}{List-II}\\\hline
(I)&\bfigi{92.3}{TOAS18_p1_phy_s2_q10_1.png}\par Transverse travelling wave moving in positive $x$-direction with speed 20 m/s&(P)&$v_p=\dfrac{3\pi}{50}$ m/s\\\hline
(II)&\bfigi{87.6}{TOAS18_p1_phy_s2_q10_2.png}\par A transverse standing wave oscillating with frequency 50 Hz, having maximum amplitude 2.5 mm&(Q)&$v_p=\dfrac\pi{10}$ m/s\\\hline
(III)&\bfigi{86.9}{TOAS18_p1_phy_s2_q10_3.png}\par A longitudinal wave propagating in negative $x$-direction with speed 20 m/s&(R)&moving in positive $x$-direction\\\hline
(IV)&\bfigi{97.5}{TOAS18_p1_phy_s2_q10_4.png}\par A longitudinal standing wave oscillating with frequency 50 Hz, having maximum amplitude 2.5 mm.&(S)&moving in positive $y$-direction\\\hline
&&(T)&Insufficient information to find direction of motion\\\hline
\end{tabular}\end{center}
\begin{qoptions}
\opt{A}{I $\to$ Q,S;II $\to$ P,T;III $\to$ Q;IV $\to$ P,T}
\opt{B}{I $\to$ Q,S;II $\to$ P;III $\to$ Q,R;IV $\to$ P,R}
\opt{C}{I $\to$ Q,S;II $\to$ P,T;III $\to$ Q,T;IV $\to$ P}
\opt{D}{I $\to$ P,T;II $\to$ Q,S;III $\to$ Q;IV $\to$ Q,T}
\end{qoptions}
\end{bankq}
\end{document}
